\documentclass[10pt,a4paper]{amsart}
\usepackage[margin=19mm]{geometry}
\usepackage{amsmath,amssymb,mathtools}
\usepackage[scr=rsfs,cal=euler]{mathalfa}
\usepackage{xfrac}
\usepackage[unicode,hidelinks,bookmarks=false]{hyperref}
\newcommand{\ie}{i.e.\ }
\newcommand{\cf}{cf.\ }
\newcommand{\resp}{respectively\ }
\newcommand{\Subsection}{\subsection}
\providecommand{\nobreakdash}{\nobreak\hbox{-}}
\setlength{\emergencystretch}{2em}
\multlinegap0pt
\usepackage{amssymb}
\usepackage{enumerate}
\newtheorem{conjecture}{Conjecture}
\newtheorem{theorem}{Theorem}
\newtheorem{proposition}{Proposition}
\def\C{\mathbb C}
\newcommand{\End}{\mathrm{End}}
\newcommand{\F}{\mathbb{F}}
\newcommand{\Ker}{\mathrm{Ker}}
\def\Q{\mathbb Q}
\def\Z{\mathbb Z}
\newcommand{\dR}{\mathrm{dR}}
\newcommand{\rk}{\mathrm{rk}}
\title{The Zilber--Pink conjecture for products of curves with highly degenerate reduction}
\author{Netan Dogra}
\date{Source: arXiv:2403.05481v3 (CC BY 4.0)}
\begin{document}
\maketitle

\noindent\emph{Source and licence.} The Zilber--Pink conjecture for products of curves with highly degenerate reduction. Version arXiv:2403.05481v3 (CC BY 4.0), distributed by arXiv under the Creative Commons Attribution 4.0 International licence, http://creativecommons.org/licenses/by/4.0/, as recorded in the arXiv licence metadata for this exact version and shown on the abstract page https://arxiv.org/abs/2403.05481v3. Full licence notice: \texttt{licenses/arxiv-2403.05481-CC-BY-4.0.txt}.\par\smallskip

\noindent\textit{Editorial scope.} Counted unit: the article's proposition prop1 (source0.tex lines 234--240). The article's complete original proof of it is delivered with it (source0.tex lines 241--251, one \texttt{proof} environment of 11 lines). No mathematics is written, repaired or re-derived: the statement and the proof are the article's own lines, in the article's own notation, and no retained line is substituted or re-flowed.  Retained uncounted as the source-local context the proof needs: lines 96--98; lines 103--114; lines 118--124.  Nothing was omitted from any retained span, so the package carries no omission record.  The retained lines cite 4 works, whose entries are transcribed verbatim from the article's own bibliography source in the e-print and delivered with short editorial labels recorded in the item's reference mapping.  No retained line contains a figure environment, an image-inclusion command or drawing code, so no image is delivered and the package declares no asset dependency.  Editorial formatting only, in addition: an AMS \texttt{amsart} preamble that restates the article's own declarations and macros used by the retained lines, namely the environments \texttt{conjecture}, \texttt{theorem}, \texttt{proposition} on separate counters, together with the standard packages those lines need.  The retained environments print in this document as Conjecture 1, Theorem 1, Proposition 1 in order of appearance, whereas the article prints the same items with the numbering of its own full document.  The complete native source of the article as distributed in the arXiv e-print for this exact version is bundled unchanged as \texttt{sources/155794/source0.tex}.
\begin{equation}\label{eqn:ker}
\Ker (\sum m_i \pi _{i *}),
\end{equation}

\begin{conjecture}\label{conj1}
Let $X$ be a smooth projective genus $g$ curve over $\mathbb{C}$, and $n>0$. Then $X^n (\mathbb{C})_{\rk <n}$ is not Zariski dense in $X^n$.
\end{conjecture}
The general form of the Zilber--Pink conjecture for abelian varieties is as follows (note that we will say nothing about this general form in the remainder of the paper).
\begin{conjecture}\label{conj2}
Let $A$ be an abelian variety of dimension $g$ over $\mathbb{C}$. Let $V\subset A$ be an irreducible subvariety of dimension $d$ not contained in a proper subgroup of $A$. Then 
\[
V(\C )\cap \cup _{B\subset A}B(\C )
\]
is not Zariski dense in $V$, where the union is over all algebraic subgroups $B\subset A$ of codimension at least $d+1$.
\end{conjecture}
If $\End (J)=\Z $, then Conjecture \ref{conj1} is equal to Conjecture \ref{conj2} with $(A,V)=(J^n ,X^n )$, since all subgroups of $J^n$ will be of the form \eqref{eqn:ker}. Otherwise conjecture 1 is strictly weaker, as it does not involve the subgroups of $J^n$ arising from extra endomorphisms. Cases of conjecture 2 when $A$ has complex multiplication are proved by R\'emond in \cite{rem3}, and when $A$ is a product of elliptic curves by Hubschmid and Viada \cite{HV}. In \cite{BD22}, Barroero and Dill showed that to prove Conjecture \ref{conj2}, it is enough to prove it over number fields. In \cite{HP}, Habegger and Pila show how, using the Pila--Zanier strategy, the conjecture would follow from suitable Galois lower bounds on `optimal singletons'.

\begin{theorem}\label{thm:main1}
Let $X/\overline{\Z }_p $ be a stable genus $g$ curve with smooth generic fibre $X_{\overline{\Q }_p }$. Let $\Gamma $ denote the dual graph of $X_{\overline{\F }_p }$ in the sense above. Suppose that
\[
n\leq \min \{g-\# E(S)-2\sum _{v\in S}^n g(v ):S\subset V(\Gamma ), \# S\leq n \}.
\]
Then Conjecture \ref{conj1} holds for $X_{\mathbb{C}}$ and $n$.
\end{theorem}

\begin{proposition}\label{prop1}
For $X$ as in the statement of Theorem \ref{thm:main1}, for each $n$-tuple $(v_1 ,\ldots ,v_n )$ of elements of $V(\Gamma ) \sqcup E(\Gamma )$, there exists a subspace $W$ of $H^0 (X,\Omega )$ of dimension at least $n$ such that
\begin{enumerate}
\item $X(\overline{\Q }_p )^n _{\rk <n} \cap U_{v_1 }^+ \times \ldots \times U_{v_n }^+ $ maps to $Z_n (W^* )$ under the map $\log _{W}^n $.
\item For all $\omega \in W$, the restriction of $\int \omega $ to $U_v ^+ $ is a rigid analytic function on $U_v ^+ $.
\end{enumerate}
\end{proposition}

\begin{proof}
For part (1), the projection map
\[
H^0 (X,\Omega )^* \to W^*
\]
sends $Z_n (H^0 (X,\Omega )^* )$ to $Z_n (W^* )$. For part (2), by the above if $\omega \in H^0 (X,\Omega )$ and $v\in V(\Gamma )\sqcup E(\Gamma )$ then the restriction of $\int \omega $ to $U_v ^+$ is a rigid analytic function if and only if the image of $\omega $ in $H^1 _{\dR}(U_v ^+ /\overline{\Q }_p )$ is zero. Let 
\[
W:=\Ker (H^0 (X,\Omega )\to \oplus _{i=1}^n H^1 _{\dR}(U_{v_i }/\overline{\Q }_p )).
\]
The conditions of Theorem \ref{thm:main1} now guarantee that $\dim (W)\geq n$.
\end{proof}

\begin{thebibliography}{99}
\bibitem[BD22]{BD22} Fabrizio Barroero and Gabriel~A. Dill. \newblock On the {Z}ilber-{P}ink conjecture for complex abelian varieties. \newblock {\em Ann. Sci. \'{E}c. Norm. Sup\'{e}r. (4)}, 55(1):261--282, 2022.
\bibitem[HP]{HP} Philipp Habegger and Jonathan Pila. \newblock O-minimality and certain atypical intersections. \newblock {\em Ann. Sci. \'{E}c. Norm. Sup\'{e}r. (4)}, 49(4):813--858, 2016.
\bibitem[HV]{HV} Patrik Hubschmid and Evelina Viada. \newblock An addendum to the elliptic torsion anomalous conjecture in codimension 2. \newblock {\em Rend. Semin. Mat. Univ. Padova}, 141:209--220, 2019.
\bibitem[rem3]{rem3} Ga\"{e}l R\'{e}mond. \newblock Intersection de sous-groupes et de sous-vari\'{e}t\'{e}s. {III}. \newblock {\em Comment. Math. Helv.}, 84(4):835--863, 2009.
\end{thebibliography}
\end{document}
